%=====================================================================================================================
\section{Two levers}\label{sec:levers}
Equation~\ref{eq:law} suggests two ways to keep a long rollout near its opening without retraining: raise $\alpha$, or
move $\mu$ toward the opening.

\paragraph{Raising $\alpha$: a content-aware correction.} At denoising step 1 of 4, the noisy chunk is denoised once
more with a temporary cache that holds only the anchor (the opening chunk). The difference between the per-channel
spatial means of the two predictions (the DC tone) is added to the real prediction, scaled by $\beta$; the main cache
and the random stream are untouched. In a preregistered confirmation on 40 held-out prompts $\times$ 4 seeds
(study~2), colour drift $D$ fell by 1.48 \ci{1.19}{1.76} at $\beta=1$ and by 1.78 \ci{1.45}{2.11} at $\beta=1.5$, with
structure drift $K$ non-inferior. Table~\ref{tab:baselines} compares it with post-hoc colour matching to the anchor,
which reaches a similar $D$ but, by construction, removes legitimate global changes too: on 7 held-out prompts with
real brightness transitions (tunnels, garages, clouds), matching keeps 7--12\% of the transition and the correction a
median of 83\% \ci{67}{147} ($\beta=1$) and 93\% \ci{83}{148} ($\beta=1.5$; prompt-bootstrap intervals over only
seven prompts). The per-prompt values range from 0.58 to
1.71, so on some prompts the correction amplifies the change. The costs are one extra forward pass per chunk (about
25\% more compute), more frame-to-frame jitter ($\times$1.13 / $\times$1.67) and larger jumps at chunk boundaries
($\times$1.32 / $\times$2.40).

A deeper attention sink, as in Deep Forcing~\cite{deepforcing}, is a strong training-free alternative (study 13, same
prompts and seeds, bit-identical pairing). Keeping 6 or 9 latent frames as the sink cuts colour drift by 1.30 and 1.93,
the latter more than the correction at $\beta=1$ ($-0.45$ \ci{-0.75}{-0.15}), with frame-to-frame jitter similar to the
correction's ($\times$1.04, $\times$1.16) and smaller boundary jumps ($\times$0.98, $\times$1.09). A sink pulls everything
toward the opening through attention, and in two of the seven transition prompts the nine-frame sink erased most of a
legitimate change (7\% and 19\% kept, against 151\% and 147\% for the correction); its median is 0.59 \ci{0.19}{1.09}
against 0.83 for the correction, but with seven prompts the intervals overlap. The correction, which moves only the DC
tone, is the more conservative choice when real changes matter.

The denoising step is a knob. On the development set, correcting at step 0 removes as much drift as step 1 but erases
tunnels (56\% of transitions kept vs 94\%); steps 2--3 do nothing, consistent with global tone being fixed at the
earliest steps~\cite{lin2024,offsetnoise}. In the relaxation study the correction raised Self Forcing's persistence from
0.28 to 0.40 ($+0.12$ \ci{0.05}{0.20}, preregistered) and also the zero-drift level by 5~\Lstar.

\begin{table}[t]
\centering\small
\caption{Correction versus post-hoc colour matching and deeper sinks (study 2 confirmation set: 40 held-out prompts
$\times$ 4 seeds; transitions: 7 held-out prompts with legitimate brightness changes). Transitions kept: median over
prompts of the corrected over the uncorrected first-to-last brightness change, prompt-bootstrap 95\% CI.}
\label{tab:baselines}
\begin{tabular}{@{}lcccc@{}}
\toprule
Method & $\Delta D$ [95\% CI] & transitions kept & jitter & boundary jumps \\
\midrule
Content-aware correction, $\beta=1$ & $-1.48$ \ci{-1.76}{-1.19} & 0.83 \ci{0.67}{1.47} & $\times$1.13 & $\times$1.32 \\
Content-aware correction, $\beta=1.5$ & $-1.78$ \ci{-2.11}{-1.45} & 0.93 \ci{0.83}{1.48} & $\times$1.67 & $\times$2.40 \\
Post-hoc mean matching & $-0.89$ \ci{-1.29}{-0.50} & 0.07 & $\times$0.15 & $\times$0.12 \\
Post-hoc Reinhard~\cite{reinhard} & $-1.58$ \ci{-1.96}{-1.19} & 0.12 & $\times$0.13 & $\times$0.10 \\
Deeper sink, 6 frames (study 13) & $-1.30$ \ci{-1.68}{-0.93} & 0.70 \ci{0.58}{1.12} & $\times$1.04 & $\times$0.98 \\
Deeper sink, 9 frames (study 13) & $-1.93$ \ci{-2.31}{-1.55} & 0.59 \ci{0.19}{1.09} & $\times$1.16 & $\times$1.09 \\
\bottomrule
\end{tabular}
\end{table}

\paragraph{Moving $\mu$: prompts.} Since the prompt sets the default, lighting words can be chosen so that the default
matches the opening. A free-form Qwen2.5-VL~\cite{qwenvl} caption of the opening does not reduce brightness drift on
either model (Self Forcing $-1.21$ \ci{-3.09}{0.36}, study 3): it calls half of the normally lit frames ``dim'', though it
improves content retention. Choosing the lighting sentence whose \emph{measured} default is nearest the opening's
brightness reduces $|\Delta L|$ at 31\,s by 26\% on Self Forcing ($-2.45$ \ci{-4.08}{-1.04}, Holm $p=0.013$, study 4);
the other five preregistered prompt tests failed, including LongLive with its own calibration (study 5). The mapping
from words to brightness is model-specific and must be measured per model. Appendix~\ref{app:default} shows how well
Eq.~\ref{eq:law} forecasts these outcomes ($r=0.82$ for the change, $0.46$ for the gain).
